\documentclass{article}
\usepackage{icmlko}
\usepackage{graphicx}
\usepackage{booktabs}
\usepackage{amsmath}

\title{「식별 등급」은 절반만 이름이었다 --- 다섯 서술 통계는 1.3842 비트 중 0.362 비트만 나른다 \\
  {\large 그리고 남은 1.0221 비트는 같은 파일의 \emph{「등급 미사용」} 칸에 통째로 적혀 있었다}}

\begin{document}
\maketitle

\begin{abstract}
직전 5개 사이클(901--905)은 \textbf{식별 등급 A/B/C} 라는 라벨을 두고 다퉜다.
프레시 티처가 물음을 하나로 좁혔다 --- \textbf{그 라벨이 실제로 나르는 정보가 몇 비트인가.}
이 노트는 그 물음에 답하면서 \textbf{물음 자체가 놓인 자리}를 먼저 고친다:
다섯 서술 통계의 조합은 T1 81짝에서 \textbf{81가지}로 서로 달라
\textbf{암기표가 곧 「재현 함수」}이고 조건부 엔트로피는 정의상 0 이 된다.
그래서 재현 여부가 아니라 \textbf{복잡도를 통제한 재현}으로 답했다.
\claimx{다섯 서술 통계는 등급의 1.3842 비트 중 0.362 비트만 나른다 --- 잔여가 1.0221 비트 있다.}
{도메인을 포함한 다섯 통계로 학습한 규칙이 leave-one-domain-out 에서 1짝도 안 틀리면 무너진다.
실측 \textbf{22} 짝을 틀린다(분모 T1 \textbf{81}). 반대로 재대입은 \textbf{0} 까지 내려가고,
그 \textbf{0} 은 자료의 성질이 아니라 튜플이 유일하다는 성질이다.}
\claimx{그 잔여는 같은 파일이 「등급 계산에 안 들어간다」고 선언한 칸에 통째로 적혀 있었다.}
{W 코드 집합이 등급을 결정하지 못하는 짝이 하나라도 있으면 무너진다.
실측 --- W 코드 집합은 \textbf{5} 가지뿐이고 등급이 섞인 칸이 \textbf{0} 이다.
상호정보 \textbf{1.3842} 비트 = 등급 엔트로피 전량.}
$\Rightarrow$ 이 노트가 뒤집는 것은 등급도 자료도 아니라 \textbf{「무엇이 이미 산출물에 적혀 있었나」를
사람이 훑는다는 절차}다. 러너 2개 절이 그것을 기계로 했다.
\end{abstract}

\section{물음을 글자대로 받으면 답이 자명하다}

티처 \#68 이 지정한 물음은 이것이다 ---
\textbf{(형 $\cdot$ 도메인 $\cdot$ 값 가짓수 $\cdot$ 비결측 $\cdot$ 소수 쪽 크기)}
다섯 서술 통계만으로 105짝의 등급을 \textbf{완전히 재현하는 함수가 있는가.}
헤드라인 분모는 T1 \textbf{81}
(\texttt{runners/out906\_grade.json:1 분모 (조항 60 --- 매 수마다 병기한다)/T1 수(분모)}),
T2 \textbf{24} 는 \texttt{docs/prereg\_901\_intervention.md:61} 대로 \textbf{안 더한다.}
등급 분포는 A \textbf{21} $\cdot$ B \textbf{47} $\cdot$ C \textbf{13} 이고
엔트로피는 \textbf{1.3842} 비트다
(\texttt{:3-가 팔 ㄱ $\cdots$/엔트로피 $\cdot$ T1 81/H(등급)}).

\textbf{그런데 그 다섯 통계의 조합은 T1 81짝에서 81가지로 서로 다르다}
(\texttt{:3-가 팔 ㄱ $\cdots$/엔트로피 $\cdot$ T1 81/5-튜플 암기/서로 다른 튜플 수} = \textbf{81} $\cdot$
유일 비율 \textbf{1.0} $\cdot$ 암기 재대입 오류 \textbf{0}).
즉 조회표를 만들면 오류가 0 이고 $H(\text{등급}\mid\text{다섯 통계})=\mathbf{0.0}$ 이 된다.
\textbf{물음을 글자대로 받으면 답은 「있다」이고, 그 답은 자료가 아니라 표본 크기가 만든다.}

그래서 이 노트는 물음을 세 겹으로 다시 놓는다:
\textbf{(1)} 규칙의 표현력을 결정트리 깊이로 제한한 \textbf{재대입},
\textbf{(2)} 도메인 하나씩 빼고 학습해 뺀 도메인에서 채점하는 \textbf{LODO},
\textbf{(3)} 조회표를 \textbf{암기 상한}으로 나란히 싣기.
사전등록이 측정 전에 이 셋을 못 박았다(\texttt{docs/prereg\_906\_grade.md} $\S$1).

\section{다섯 통계가 실제로 나르는 것 --- 0.362 비트}

\begin{figure}[t]\centering
\includegraphics[width=\linewidth]{figs/fig1.pdf}
\caption{(a) 재대입은 깊이가 늘면 0 까지 내려가지만 LODO 는 22 에서 멈춘다.
(b) 조건별 상호정보. 붉은 막대(다섯 통계 전부)는 자료가 아니라 튜플이 유일해서 생긴 값이다.
값은 전부 \texttt{runners/out906\_grade.json} 에서 읽었다 --- 손 전사 0.}
\end{figure}

\begin{table}[t]\centering\small
\begin{tabular}{lrr}
\toprule
특징 집합 & 최소 재대입 오류 & 최소 LODO 오류 \\
\midrule
F5 다섯 통계(전부) & \textbf{0} & \textbf{22} \\
F4 도메인 뺀 넷 & \textbf{0} & \textbf{22} \\
F3 소수 쪽 뺀 넷(음성 대조) & \textbf{7} & \textbf{25} \\
\bottomrule
\end{tabular}
\caption{분모는 T1 \textbf{81} 짝. 출처
\texttt{runners/out906\_grade.json:3-나 팔 ㄱ $\cdots$/$\langle$특징 집합$\rangle$/최소 재대입 오류 $\cdot$ 최소 LODO 오류}.}
\end{table}

일반화되는 규칙의 잔여는 $H(\text{등급}\mid\hat{y}_{\mathrm{LODO}}) = \mathbf{1.0221}$ 비트이고
상호정보는 \textbf{0.362} 비트다
(\texttt{:3-나 $\cdots$/잔여 비트(= H(등급 | LODO 예측))} $\cdot$ \texttt{/일반화 기준 상호정보(비트)}).
\textbf{1.3842} 비트 중 \textbf{0.362} 비트다.

\textbf{음성 대조가 중요하다.} 소수 쪽은 \texttt{식1 근거} 블록 안에 있다 ---
즉 식1 을 판정한 그 루틴의 산물이다. 그것을 뺀 F3 은 최소 재대입 오류가
\textbf{0} 에서 \textbf{7} 로, 최소 LODO 오류가 \textbf{22} 에서 \textbf{25} 로 나빠질 뿐이다.
\textbf{답을 물음에서 베끼는 자리를 막아도 그림이 안 바뀐다}(이슈 \#155 처방 2).

\textbf{그리고 도메인은 재대입에서 아무 일도 안 한다} --- F5 와 F4 의 최소 재대입 오류가
둘 다 \textbf{0}, 최소 LODO 오류가 둘 다 \textbf{22} 다.

\subsection{잔여를 지는 짝은 어디에 있나}

LODO 가 틀리는 \textbf{22} 짝(\texttt{:3-나 $\cdots$/잔여를 지는 짝 수})은
도메인 「팝업」에 \textbf{15} 짝이 몰려 있고,
식3 으로 가르면 예 \textbf{14} $\cdot$ 모른다 \textbf{8} $\cdot$ 아니오 \textbf{3} 이다.
\textbf{그 15짝은 전부 \texttt{intervention.attributes.*} 접두사를 갖는다} ---
잔여를 가르는 것은 수치가 아니라 \textbf{열 이름의 뜻}이라는 뜻이고,
\textbf{저장소에는 열 이름의 뜻을 담은 사전이 없다}(사전등록 $\S$⓪-나).

\section{그런데 그 잔여는 같은 파일에 이미 적혀 있었다}

\begin{figure}[t]\centering
\includegraphics[width=\linewidth]{figs/fig2.pdf}
\caption{(a) W 코드 집합 5가지가 등급을 결정한다 --- 섞인 칸 0.
(b) 잔여 22짝의 도메인 분해.}
\end{figure}

이 노트가 스스로 고른 팔은 이것이다 ---
\textbf{같은 파일이 「등급 계산에 안 들어간다」고 선언한 형제 키들이 정보 차원에서도 그런가.}

\texttt{runners/out902\_identify.json} 은 최상위 키
\textbf{「W 는 어느 것도 등급 계산에 안 들어간다」} 를 갖고 있다.
그 W 목록을 코드 집합으로 읽으면 T1 81짝에서 \textbf{5} 가지뿐이고
(\texttt{runners/out906\_grade.json:3-다 팔 ㄴ $\cdots$/W 코드 집합만/특징 값 가짓수}),
\textbf{등급이 2가지 이상 섞인 칸이 0} 이다
(\texttt{:$\cdots$/대응이 일대일인가/등급이 두 가지 이상 섞인 W 코드 집합 수}).
상호정보는 \textbf{1.3842} 비트, 즉 \textbf{등급 엔트로피 전량}이고
조건부 엔트로피는 \textbf{0.0} 이다.

\begin{table}[t]\centering\small
\begin{tabular}{lrrr}
\toprule
W 코드 집합 & A & B & C \\
\midrule
W2, W8 & \textbf{21} & & \\
W2, W4, W8 & & \textbf{45} & \\
W2, W6, W8 & & 2 & \\
W2, W7, W8 & & & 3 \\
W2, W4, W6, W8 & & & \textbf{10} \\
\bottomrule
\end{tabular}
\caption{분모 T1 \textbf{81} 짝. 출처
\texttt{runners/out906\_grade.json:3-다 팔 ㄴ $\cdots$/W 코드 집합 → 등급 대응표(T1 81)}.}
\end{table}

\textbf{놀랍지 않은 부분과 놀라운 부분을 갈라 적는다.}
안 놀라운 것 --- W 는 「무엇이 없어서 못 재는가」의 목록이라 조건에서 파생되고,
그러니 등급을 결정하는 것 자체는 대수롭지 않다.
\textbf{놀라운 것} --- 그 파생물이 \textbf{「등급 미사용」이라는 이름표를 달고} 산출물에 실려 있었고,
\textbf{「등급이 나르는 잔여가 무엇인가」를 5개 사이클 동안 물으면서 아무도 그 이름표를 안 열었다.}
그리고 그 잔여(W4 $\cdot$ W6 $\cdot$ W7 의 유무)는
\textbf{다섯 서술 통계가 못 나르는 바로 그 부분이다.}

\section{장치 --- 자백을 기계가 깎는다}

티처 \#68 은 직전 노트의 자백 장치를 2군데서 부쉈다:
P2 가 \textbf{답을 물음에서 베끼는 항등식}이었고(C1),
$\S$0 과 예측의 교집합 0 은 \textbf{저자가 지은 두 목록의 차이}였다(C2).
이슈 \#155 의 처방을 이번엔 글자가 아니라 \textbf{기계}로 실행했다.

\claimx{예측마다 그 예측을 거짓으로 만드는 입력을 심었고, 장치가 반증 불가능한 것을 실제로 뺐다.}
{심은 반증 입력에서 평가기가 「확인」을 그대로 내면 그 예측은 명부에서 자동으로 빠진다.
실측 --- 심은 \textbf{4} 중 반증 가능으로 확인된 것 \textbf{3},
명부에 남은 예측 \textbf{3}(확인 \textbf{2} $\cdot$ 반증 \textbf{1}),
장치가 뺀 것 \textbf{1}(반증 불가능하게 설계한 가짜 예측).}

배선 검사는 심은 결함 \textbf{11} 에 대해 발화 \textbf{11} $\cdot$ 국소 \textbf{11} $\cdot$
음성 대조 \textbf{11} $\cdot$ 지우면 통과 \textbf{11} $\cdot$ 계수기 증가 \textbf{11} 이고,
직전 노트에서 \textbf{검정력이 0} 이라 지적받은 「국소」 열에 대해
\textbf{그 열이 거짓이 되는 입력}을 \textbf{8} 개 심어 \textbf{8} 개 전부 발화시켰다
(\texttt{runners/out906\_grade.json:2 배선 검사/㉱ 국소 열 검정력 $>$ 0 인 검사 수}).

그리고 티처 \#68 ⑧ 의 처방 --- \textbf{「지금 읽고 있는 그 산출물의 최상위 키를 전량 읽고,
오늘 낼 문장이 그중 어디에 이미 있는지 대조하라」} --- 를 \textbf{러너 2개 절}로 넣었다.
입력 산출물의 최상위 키 \textbf{30} 개를 전량 찍고,
판정문에서 뽑은 정수 \textbf{9} 가지를 입력에서 되찾아
\textbf{6} 가지는 이미 있고 \textbf{3} 가지(엔트로피 값들)만 오늘 새것임을 기계로 적었다
(\texttt{:0-나 $\cdots$/입력에 이미 있는 수의 개수} $\cdot$ \texttt{/입력에서 못 찾은 수의 개수}).
\textbf{양쪽이 다 기계 산물이다} --- 판정문은 러너가 만들고 경로는 탐색이 찾는다.

\section{그래서 무엇이 정해졌나}

\textbf{등급은 절반만 이름이었다.} 다섯 서술 통계가 \textbf{0.362} 비트를 나르고
\textbf{1.0221} 비트가 남는다. 남는 쪽은 식3(역인과 아님)과 식1 의 표본 크기 미달이고,
그 둘은 \textbf{같은 파일의 W 목록에 이미 코드로 적혀 있다.}
그러므로 다음 수는 \textbf{등급을 더 다투는 것이 아니라}
\textbf{잔여를 지는 22짝의 열 이름이 무엇을 뜻하는지 기계가 읽을 수 있게 만드는 것}이다.
그 사전은 저장소에도 세상에도 없다 --- 우리 자료의 우리 열이라 밖에서 받을 것이 없다.

\end{document}
